\chapter{Final Exam Summary: Formulae and Nuances}

\section{Golden rules (read these first)}

These nine rules cause most of the lost marks in exams. Learn them before the formulas.

\begin{enumerate}
  \item \textbf{Count the tails.}
  \begin{itemize}
    \item Two-sided: split $\alpha$ into two parts, $\alpha/2$ in each tail. Use $z_{\alpha/2}$, $t_{\alpha/2}$, $\chi^2_{\alpha/2}$ and $\chi^2_{1-\alpha/2}$, $f_{\alpha/2}$ and $f_{1-\alpha/2}$.
    \item One-sided: put all of $\alpha$ in one tail. Use $z_\alpha$, $t_\alpha$, $\chi^2_\alpha$ or $\chi^2_{1-\alpha}$, $f_\alpha$ or $f_{1-\alpha}$.
    \item A one-sided bound has one limit only. The other side is open.
  \end{itemize}
  \item \textbf{A one-sided $95\%$ bound is not the end of a two-sided $95\%$ interval.}
  \begin{itemize}
    \item The end of a two-sided $95\%$ interval is a one-sided $97.5\%$ bound.
    \item A one-sided $95\%$ bound is the end of a two-sided $90\%$ interval.
    \item \emph{Reason.} Every $100(1-\alpha)\%$ statement allows a total error of $\alpha$. A two-sided interval can be wrong in two ways ($\mu$ above $U$ or $\mu$ below $L$), so $\alpha$ is shared: $2.5\%$ per side for $95\%$. A one-sided bound can be wrong in one way only, so the whole $5\%$ is used on that side.
    \item \emph{Multipliers.} $1.960$ belongs to two-sided $95\%$ (each end) and to one-sided $97.5\%$. $1.645$ belongs to two-sided $90\%$ (each end) and to one-sided $95\%$.
    \item \emph{Conversion.} Each end of a two-sided interval at level $C$ is a one-sided bound at level $C+\tfrac{100-C}{2}$ (for example $95\to97.5$). A one-sided bound at level $C$ is the end of a two-sided interval at level $2C-100$ (for example $95\to90$).
    \item \emph{Numerical check} (impact-energy data: $\bar{x}=64.46$, standard error $0.316$). Two-sided $95\%$: $(63.84,\,65.08)$. One-sided $95\%$: $\mu\le64.98$ or $\mu\ge63.94$. Two-sided $90\%$: $(63.94,\,64.98)$, which equals the two one-sided $95\%$ bounds. One-sided $97.5\%$ upper bound: $65.08$.
  \end{itemize}
  \item \textbf{Match each test with the right interval, using the same $\alpha$.}
  \begin{itemize}
    \item Two-sided test at $\alpha$ $\leftrightarrow$ two-sided $100(1-\alpha)\%$ interval.
    \item One-sided test at $\alpha$ $\leftrightarrow$ one-sided $100(1-\alpha)\%$ bound.
    \item $H_1:\theta>\theta_0$ goes with the lower bound: reject if $\theta_0$ is below it. $H_1:\theta<\theta_0$ goes with the upper bound.
    \item A one-sided test (cutoff $1.645$) can reject while $\theta_0$ is still inside the two-sided $95\%$ interval (cutoff $1.96$). This is not a contradiction. Compare a one-sided test with a one-sided bound, or with the two-sided $90\%$ interval.
  \end{itemize}
  \item \textbf{Put the claim you want to prove in $H_1$.} Rejecting $H_0$ is a strong conclusion. Failing to reject only means there is not enough evidence. It never proves $H_0$.
  \item \textbf{Read the $P$-value correctly.} $P$ is the probability of a result at least this extreme if $H_0$ is true. It is not the probability that $H_0$ is true. Reject $H_0$ when $P\le\alpha$.
  \item \textbf{Read confidence correctly.} The $95\%$ describes the method: $95\%$ of intervals made this way contain the parameter. It does not give a probability for one computed interval.
  \item \textbf{Know the three intervals.} A confidence interval bounds a parameter. A prediction interval bounds one future value and is always wider. A tolerance interval bounds a stated proportion of the population.
  \item \textbf{State and check the assumptions.} Normal population for $t$, $\chi^2$ and $F$. Independent samples. Equal variances for the pooled $t$-test and for ANOVA.
  \item \textbf{Use the eight steps for every test.} Parameter; $H_0$; $H_1$; $\alpha$; test statistic; rejection rule; computation; conclusion in words.
\end{enumerate}

Common $z$ values:

\begin{center}
\small
\begin{tabular}{lcc}
\toprule
Confidence & two-sided: $z_{\alpha/2}$ & one-sided: $z_\alpha$ \\
\midrule
$90\%$ & $1.645$ & $1.282$ \\
$95\%$ & $1.960$ & $1.645$ \\
$99\%$ & $2.576$ & $2.326$ \\
\bottomrule
\end{tabular}
\end{center}

Full tables for $z$, $t$, $\chi^2$ and $F$, with one-tail and two-tail columns, are in Chapter~16.

\section{Chapter 7: Point Estimation (book Chapter~1)}

\begin{formulabox}{Judging an estimator}
\begin{itemize}
  \item \textbf{Unbiased; bias:} $\Exp[\hat\Theta]=\theta$
  \item $\text{bias}=\Exp[\hat\Theta]-\theta$
  \item \textbf{Standard error:} $\se(\hat\Theta)=\sqrt{\Var[\hat\Theta]}$
\end{itemize}
\[ \se(\bar{X})=\frac{\sigma}{\sqrt n} \]
\begin{itemize}
  \item \textbf{Mean square error:} $\MSE(\hat\Theta)=\Var[\hat\Theta]+(\text{bias})^2$
  \item \textbf{Relative efficiency (less than 1 means $\hat\Theta_1$ is better):} $\MSE(\hat\Theta_1)\big/\MSE(\hat\Theta_2)$
\end{itemize}
\end{formulabox}

\begin{formulabox}{Building estimators}
\textbf{Method of moments: equate $\Exp[X^k]$ to $\frac1n\sum X_i^k$ and solve}
\[ \text{Exponential: }\hat\lambda=\frac1{\bar{x}} \]
\begin{itemize}
  \item $\text{Uniform}(0,a):\ \hat{a}=2\bar{x}$
  \item \textbf{Normal (moments and likelihood agree):} $\hat\mu=\bar{x}$
\end{itemize}
\[ \hat\sigma^2=\frac1n\sum(x_i-\bar{x})^2 \]
\textbf{Gamma by moments, with $m_2=\frac1n\sum x_i^2$}
\[ \hat{r}=\frac{\bar{x}^2}{m_2-\bar{x}^2} \]
\[ \hat\lambda=\frac{\bar{x}}{m_2-\bar{x}^2} \]
\begin{itemize}
  \item \textbf{Maximum likelihood: maximise $\prod f(x_i;\theta)$, usually through $\ln L$:} $\text{Bernoulli: }\hat{p}=\bar{x}$
\end{itemize}
\[ \text{Exponential: }\hat\lambda=\frac1{\bar{x}} \]
\begin{itemize}
  \item $\text{Uniform}(0,a):\ \hat{a}=\max x_i$
\end{itemize}
\end{formulabox}

\begin{formulabox}{Sampling distributions}
\textbf{One mean}
\[ \bar{X}\sim\Normal\!\Big(\mu,\frac{\sigma^2}{n}\Big) \]
\textbf{Difference of two independent means}
\[ \begin{aligned}\bar{X}_1-\bar{X}_2\ &\sim\ \Normal(\mu_1-\mu_2,\ \sigma_D^2)\\ \sigma_D^2\ &=\ \frac{\sigma_1^2}{n_1}+\frac{\sigma_2^2}{n_2}\end{aligned} \]
Exact for normal populations; approximate by the CLT when $n\ge30$.
\end{formulabox}

\begin{notebox}
\begin{itemize}
  \item $\bar{X}$, $S^2$ (divisor $n-1$) and $\hat{P}=X/n$ are unbiased, but $S$ is biased for $\sigma$ ($\Exp[S]=c_4\sigma$; $c_4=0.94,\ 0.9727,\ 0.9896$ for $n=5,10,25$).
  \item A biased estimator can have a smaller MSE than an unbiased one. $\bar{X}$ is the MVUE of $\mu$ for a normal sample.
  \item MLE \emph{invariance}: the MLE of $h(\theta)$ is $h(\hat\theta)$, so the MLE of $\sigma$ is $\sqrt{\hat\sigma^2}$, not $s$. The normal $\hat\sigma^2$ has divisor $n$: biased, but asymptotically unbiased.
  \item Uniform$(0,a)$: the likelihood has no stationary point, so calculus fails and $\hat{a}=\max x_i$.
  \item Bayes: posterior $\propto$ likelihood $\times$ prior; Bernoulli with a $\mathrm{Beta}(a,b)$ prior gives $\mathrm{Beta}(a+x,\,b+n-x)$. Bootstrap: resample with replacement; the spread of $\hat\theta^*$ estimates $\se(\hat\theta)$.
\end{itemize}
\end{notebox}

\subsubsection*{Worked numericals}

\begin{example}[Maximum likelihood for an exponential rate]
\begin{itemize}
  \item \textbf{Given.} Eight failure times with mean $\bar{x}=21.65$ hours. Model: exponential, $\Exp[X]=1/\lambda$.
  \item \textbf{Estimator.} $\hat\lambda=1/\bar{x}$ (method of moments and maximum likelihood agree).
  \item \textbf{Compute.} $\hat\lambda=\dfrac{1}{21.65}=0.0462$ per hour.
  \item \textbf{Conclusion.} The estimated failure rate is $0.0462$ per hour.
\end{itemize}
\end{example}

\begin{example}[Comparing two estimators]
\begin{itemize}
  \item \textbf{Given.} Estimator 1: unbiased, variance $4$. Estimator 2: bias $1$, variance $2$.
  \item \textbf{MSE of 1.} $\MSE_1=4+0^2=4$.
  \item \textbf{MSE of 2.} $\MSE_2=2+1^2=3$.
  \item \textbf{Relative efficiency.} $\MSE_2/\MSE_1=3/4=0.75<1$.
  \item \textbf{Conclusion.} Estimator 2 is better in the MSE sense even though it is biased.
\end{itemize}
\end{example}

\begin{example}[Probability about a sample mean]
\begin{itemize}
  \item \textbf{Given.} Normal population, $\sigma=10$, $n=25$, $\mu=100$. Find $\Prob(\bar{X}<95)$.
  \item \textbf{Standard error.} $\se(\bar{X})=10/\sqrt{25}=2$.
  \item \textbf{Standardise.} $z=\dfrac{95-100}{2}=-2.5$.
  \item \textbf{Probability.} $\Prob(Z<-2.5)=0.0062$.
  \item \textbf{Conclusion.} There is about a $0.6\%$ chance that the sample mean is below $95$.
\end{itemize}
\end{example}

\begin{example}[Difference of two sample means]
\begin{itemize}
  \item \textbf{Given.} Old process: $\mu=5000$, $\sigma=40$, $n=16$. New process: $\mu=5050$, $\sigma=30$, $n=25$. Find $\Prob(\bar{X}_{\text{new}}-\bar{X}_{\text{old}}\ge25)$.
  \item \textbf{Mean.} $50$.
  \item \textbf{Variance.} $\dfrac{30^2}{25}+\dfrac{40^2}{16}=36+100=136$, so the standard deviation is $11.66$.
  \item \textbf{Standardise.} $z=\dfrac{25-50}{11.66}=-2.14$.
  \item \textbf{Probability.} $\Prob(Z\ge-2.14)=0.9838$.
  \item \textbf{Conclusion.} The new process beats the old by at least $25$ hours with probability $0.984$.
\end{itemize}
\end{example}

\section{Chapter 8: Statistical Intervals (book Chapter~2)}

\begin{formulabox}{Two-sided $100(1-\alpha)\%$ intervals}
\textbf{Mean, $\sigma$ known}
\[ \bar{x}\pm z_{\alpha/2}\,\frac{\sigma}{\sqrt n} \]
\textbf{Mean, $\sigma$ unknown (for $n\ge40$ use $z_{\alpha/2}$)}
\[ \bar{x}\pm t_{\alpha/2,\,n-1}\,\frac{s}{\sqrt n} \]
\textbf{Variance}
\[ \begin{aligned}\sigma^2&\ \ge\ \frac{(n-1)s^2}{\chi^2_{\alpha/2,\,n-1}}\\[4pt] \sigma^2&\ \le\ \frac{(n-1)s^2}{\chi^2_{1-\alpha/2,\,n-1}}\end{aligned} \]
\textbf{Proportion (large $n$)}
\[ \hat{p}\pm z_{\alpha/2}\sqrt{\frac{\hat{p}(1-\hat{p})}{n}} \]
\textbf{Prediction interval for one new value}
\[ \bar{x}\pm t_{\alpha/2,\,n-1}\,s\sqrt{1+\frac1n} \]
\begin{itemize}
  \item \textbf{Tolerance interval ($k$ from Table XI):} $\bar{x}\pm ks$
\end{itemize}
\end{formulabox}

\begin{formulabox}{One-sided bounds: use $\alpha$ and keep one limit}
\textbf{Mean, $\sigma$ known (unknown $\sigma$: use $t_{\alpha,n-1}$)}
\[ \mu\le\bar{x}+z_\alpha\frac{\sigma}{\sqrt n} \]
\[ \mu\ge\bar{x}-z_\alpha\frac{\sigma}{\sqrt n} \]
\textbf{Variance}
\[ \sigma^2\le\frac{(n-1)s^2}{\chi^2_{1-\alpha,\,n-1}} \]
\[ \sigma^2\ge\frac{(n-1)s^2}{\chi^2_{\alpha,\,n-1}} \]
\begin{itemize}
  \item \textbf{Proportion:} $p\le\hat{p}+z_\alpha\sqrt{\tfrac{\hat{p}(1-\hat{p})}{n}}$
  \item $p\ge\hat{p}-z_\alpha\sqrt{\tfrac{\hat{p}(1-\hat{p})}{n}}$
\end{itemize}
\end{formulabox}

\begin{formulabox}{Sample size for error at most $E$ (round up)}
\textbf{Mean}
\[ n=\Big(\frac{z_{\alpha/2}\,\sigma}{E}\Big)^2 \]
\textbf{Proportion (worst case $p=0.5$)}
\[ n=\Big(\frac{z_{\alpha/2}}{E}\Big)^2p(1-p) \]
Use $z_\alpha$ for a one-sided bound. The interval length is $2E$.
\end{formulabox}

\begin{notebox}
\begin{itemize}
  \item Halving $E$ quadruples $n$. Higher confidence gives a longer interval; larger $n$ a shorter one.
  \item The $\sigma^2$ interval is not symmetric and is very sensitive to non-normality. The proportion interval needs $np\ge5$ and $n(1-p)\ge5$.
  \item Pivot method: find $g(X;\theta)$ whose distribution is free of $\theta$ and solve $a\le g\le b$ for $\theta$.
  \item Textbook values: impact energy ($\bar{x}=64.46$, $\sigma=1$, $n=10$) gives $(63.84,\,65.08)$; adhesion ($n=22$, $\bar{x}=13.71$, $s=3.55$) gives $(12.14,\,15.28)$ with $t_{0.025,21}=2.080$.
\end{itemize}
\end{notebox}

\subsubsection*{Worked numericals}

\begin{example}[Two-sided $t$ interval]
\begin{itemize}
  \item \textbf{Given.} $n=25$, $\bar{x}=50$, $s=4$, confidence $95\%$.
  \item \textbf{Critical value.} $24$ df, so $t_{0.025,24}=2.064$.
  \item \textbf{Standard error.} $s/\sqrt n=4/5=0.8$.
  \item \textbf{Interval.} $50\pm2.064(0.8)=50\pm1.65$.
  \item \textbf{Conclusion.} $48.35\le\mu\le51.65$.
\end{itemize}
\end{example}

\begin{example}[One-sided bound, same data]
\begin{itemize}
  \item \textbf{Given.} Same $n=25$, $\bar{x}=50$, $s=4$. Find the $95\%$ upper bound.
  \item \textbf{Critical value.} Use the one-tail value $t_{0.05,24}=1.711$.
  \item \textbf{Bound.} $\mu\le50+1.711(0.8)=51.37$.
  \item \textbf{Conclusion.} The upper bound $51.37$ is smaller than the two-sided upper limit $51.65$, as expected.
\end{itemize}
\end{example}

\begin{example}[Sample size for a mean]
\begin{itemize}
  \item \textbf{Given.} $\sigma=3$, error at most $E=0.5$, confidence $95\%$.
  \item \textbf{Formula.} $n=\Big(\dfrac{z_{\alpha/2}\sigma}{E}\Big)^2$.
  \item \textbf{Compute.} $n=\Big(\dfrac{1.96\times3}{0.5}\Big)^2=138.3$.
  \item \textbf{Conclusion.} Round up: $n=139$.
\end{itemize}
\end{example}

\begin{example}[Interval for a variance]
\begin{itemize}
  \item \textbf{Given.} $n=15$, $s^2=4$, confidence $95\%$ two-sided, $14$ df.
  \item \textbf{Critical values.} $\chi^2_{0.025,14}=26.119$ and $\chi^2_{0.975,14}=5.629$.
  \item \textbf{Lower limit.} $\dfrac{14(4)}{26.119}=2.14$.
  \item \textbf{Upper limit.} $\dfrac{14(4)}{5.629}=9.95$.
  \item \textbf{Conclusion.} $2.14\le\sigma^2\le9.95$. For a one-sided $95\%$ upper bound use $\chi^2_{0.95,14}=6.571$: $\sigma^2\le56/6.571=8.52$.
\end{itemize}
\end{example}

\begin{example}[Interval for a proportion]
\begin{itemize}
  \item \textbf{Given.} $40$ defectives in $n=200$, so $\hat{p}=0.20$; confidence $95\%$.
  \item \textbf{Standard error.} $\sqrt{\dfrac{0.20(0.80)}{200}}=0.0283$.
  \item \textbf{Interval.} $0.20\pm1.96(0.0283)=0.20\pm0.055$.
  \item \textbf{Conclusion.} $0.145\le p\le0.255$.
\end{itemize}
\end{example}

\section{Chapter 9: Tests on a Single Sample (book Chapter~3)}

\begin{formulabox}{Test statistics}
\textbf{Mean, $\sigma$ known}
\[ z_0=\frac{\bar{x}-\mu_0}{\sigma/\sqrt n} \]
\textbf{Mean, $\sigma$ unknown ($n-1$ df)}
\[ t_0=\frac{\bar{x}-\mu_0}{s/\sqrt n} \]
\textbf{Variance ($n-1$ df)}
\[ \chi_0^2=\frac{(n-1)s^2}{\sigma_0^2} \]
\textbf{Proportion (large $n$)}
\[ z_0=\frac{x-np_0}{\sqrt{np_0(1-p_0)}} \]
\end{formulabox}

\begin{formulabox}{Decision rule: when to reject $H_0$}
Compute the test statistic ($z_0$, $t_0$, $\chi_0^2$ or $f_0$). Find the critical value in the table. Then look at the sign in $H_1$: it tells you which tail to check.
\begin{itemize}
  \item $H_1$ has $\ne$ (two-sided): reject if the statistic is far from zero in \emph{either} direction. For $z$ and $t$, reject if $\abs{\text{statistic}}>$ critical value at $\alpha/2$.
  \item $H_1$ has $>$ (one-sided, right tail): reject if the statistic is \emph{larger than} the critical value at $\alpha$.
  \item $H_1$ has $<$ (one-sided, left tail): reject if the statistic is \emph{smaller than} minus the critical value at $\alpha$. (For $\chi^2$ and $F$, smaller than the lower-tail point.)
  \item If the statistic does not pass the critical value, do not reject $H_0$.
  \item $P$-value form of the same rule: reject if $P\le\alpha$.
\end{itemize}
Example: $z_0=2.4$. One-sided $H_1:\mu>\mu_0$ at $\alpha=0.05$: $z_\alpha=1.645$ and $2.4>1.645$, so reject. Two-sided at $\alpha=0.05$: $z_{\alpha/2}=1.96$ and $\abs{2.4}>1.96$, so reject. If $z_0=-2.4$ with $H_1:\mu>\mu_0$, the statistic is on the wrong side, so do not reject.
\end{formulabox}

\begin{formulabox}{Two-sided tests ($H_1$ uses $\ne$): reject $H_0$ if}
\textbf{$z$-test and proportion test}
\[ \abs{z_0}>z_{\alpha/2}\qquad\text{that is}\qquad z_0>z_{\alpha/2}\ \text{ or }\ z_0<-z_{\alpha/2} \]
\textbf{$t$-test}
\[ \abs{t_0}>t_{\alpha/2,\,n-1} \]
\textbf{Variance test}
\[ \chi_0^2>\chi^2_{\alpha/2,\,n-1}\qquad\text{or}\qquad\chi_0^2<\chi^2_{1-\alpha/2,\,n-1} \]
\textbf{$P$-value for the $z$-test}
\[ P=2\,[1-\Phi(\abs{z_0})] \]
Each tail holds $\alpha/2$, so look up the table at one-tail area $\alpha/2$.
\end{formulabox}

\begin{formulabox}{One-sided tests ($H_1$ uses $>$ or $<$): reject $H_0$ if}
\textbf{$z$-test and proportion test}
\[ \begin{aligned} H_1:\mu>\mu_0&:\ z_0>z_\alpha\\ H_1:\mu<\mu_0&:\ z_0<-z_\alpha\end{aligned} \]
\textbf{$t$-test}
\[ \begin{aligned} H_1:\mu>\mu_0&:\ t_0>t_{\alpha,\,n-1}\\ H_1:\mu<\mu_0&:\ t_0<-t_{\alpha,\,n-1}\end{aligned} \]
\textbf{Variance test}
\[ \begin{aligned} H_1:\sigma^2>\sigma_0^2&:\ \chi_0^2>\chi^2_{\alpha,\,n-1}\\ H_1:\sigma^2<\sigma_0^2&:\ \chi_0^2<\chi^2_{1-\alpha,\,n-1}\end{aligned} \]
\textbf{$P$-value for the $z$-test}
\[ \begin{aligned} H_1:>&:\ P=1-\Phi(z_0)\\ H_1:<&:\ P=\Phi(z_0)\end{aligned} \]
All of $\alpha$ sits in one tail, so look up the table at one-tail area $\alpha$. The sign in $H_1$ points to the tail: $>$ means right tail, $<$ means left tail.
\end{formulabox}

\begin{formulabox}{Errors and sample size ($\sigma$ known, true mean $\mu_0+\delta$)}
\begin{itemize}
  \item \textbf{Error probabilities:} $\alpha=\Prob(\text{reject }H_0\mid H_0\text{ true})$
  \item $\beta=\Prob(\text{fail to reject}\mid H_0\text{ false})$
  \item $\text{power}=1-\beta$
\end{itemize}
\textbf{$\beta$, two-sided}
\[ \begin{aligned}\beta&=\Phi\!\Big(z_{\alpha/2}-\frac{\delta\sqrt n}{\sigma}\Big)\\&\quad-\Phi\!\Big(-z_{\alpha/2}-\frac{\delta\sqrt n}{\sigma}\Big)\end{aligned} \]
\textbf{$\beta$, one-sided}
\[ \beta=\Phi\!\Big(z_\alpha-\frac{\delta\sqrt n}{\sigma}\Big) \]
\textbf{Sample size, two-sided and one-sided}
\[ n\approx\frac{(z_{\alpha/2}+z_\beta)^2\sigma^2}{\delta^2} \]
\[ n=\frac{(z_\alpha+z_\beta)^2\sigma^2}{\delta^2} \]
\end{formulabox}

\begin{formulabox}{Goodness of fit and contingency tables}
\textbf{Statistic}
\[ \chi_0^2=\sum\frac{(O_i-E_i)^2}{E_i} \]
\textbf{Degrees of freedom}
\[ \begin{aligned}\text{fit: }&\ k-p-1\\ \text{contingency: }&\ (r-1)(c-1)\\ E_{ij}&=\frac{R_iC_j}{N}\end{aligned} \]
\end{formulabox}

\begin{notebox}
\begin{itemize}
  \item A two-sided test rejects iff $\mu_0$ is outside the two-sided $100(1-\alpha)\%$ interval (duality). Reject iff $P\le\alpha$.
  \item Goodness of fit: $p$ is the number of estimated parameters; combine cells with $E_i<3$--$5$; upper tail only.
  \item Statistical significance is not practical significance: with a huge $n$, a trivial shift gives a tiny $P$.
  \item Textbook sample size: $\sigma=2$, $\delta=1$, power $0.90$ gives $n=42$. OC curves use $d=\abs{\mu-\mu_0}/\sigma$.
  \item Proportion sample size: $n\approx\Big[\dfrac{z_{\alpha/2}\sqrt{p_0(1-p_0)}+z_\beta\sqrt{p(1-p)}}{p-p_0}\Big]^2$ ($z_\alpha$ if one-sided).
\end{itemize}
\end{notebox}

\subsubsection*{Worked numericals}

\begin{example}[$z$-test (known $\sigma$)]
\begin{itemize}
  \item \textbf{Given.} $H_0:\mu=100$, $\sigma=5$, $n=25$, $\bar{x}=102.4$, $\alpha=0.05$.
  \item \textbf{Statistic.} $z_0=\dfrac{102.4-100}{5/\sqrt{25}}=2.4$.
  \item \textbf{Two-sided.} Critical value $1.96$. $\abs{2.4}>1.96$, so reject. $P=2(0.0082)=0.0164$.
  \item \textbf{One-sided ($H_1:\mu>100$).} Critical value $1.645$. $2.4>1.645$, so reject. $P=0.0082$.
  \item \textbf{Conclusion.} Both tests reject $H_0$: the mean is greater than $100$.
\end{itemize}
\end{example}

\begin{example}[$t$-test (unknown $\sigma$)]
\begin{itemize}
  \item \textbf{Given.} $H_0:\mu=50$, $n=16$, $\bar{x}=52.3$, $s=4$, $\alpha=0.05$, $15$ df.
  \item \textbf{Statistic.} $t_0=\dfrac{52.3-50}{4/\sqrt{16}}=2.3$.
  \item \textbf{Two-sided.} $t_{0.025,15}=2.131$. $2.3>2.131$, so reject.
  \item \textbf{One-sided ($H_1:\mu>50$).} $t_{0.05,15}=1.753$. $2.3>1.753$, so reject.
  \item \textbf{Conclusion.} Both tests reject $H_0$.
\end{itemize}
\end{example}

\begin{example}[Test on a variance]
\begin{itemize}
  \item \textbf{Given.} $H_0:\sigma^2=9$ against $H_1:\sigma^2>9$, $n=21$, $s^2=14$, $\alpha=0.05$, $20$ df.
  \item \textbf{Statistic.} $\chi_0^2=\dfrac{20(14)}{9}=31.11$.
  \item \textbf{Critical value.} $\chi^2_{0.05,20}=31.41$ (upper tail).
  \item \textbf{Decision.} $31.11<31.41$, so do not reject. $P\approx0.054$.
  \item \textbf{Conclusion.} The evidence is just short of significance at $\alpha=0.05$.
\end{itemize}
\end{example}

\begin{example}[Test on a proportion]
\begin{itemize}
  \item \textbf{Given.} $H_0:p=0.10$ against $H_1:p>0.10$, $n=100$, $x=16$, $\alpha=0.05$.
  \item \textbf{Statistic.} $z_0=\dfrac{16-100(0.10)}{\sqrt{100(0.10)(0.90)}}=\dfrac{6}{3}=2.0$.
  \item \textbf{Decision.} $2.0>1.645$, so reject. $P=0.0228$.
  \item \textbf{Conclusion.} The proportion is greater than $0.10$.
\end{itemize}
\end{example}

\begin{example}[Sample size for a test]
\begin{itemize}
  \item \textbf{Given.} Detect a shift $\delta=1$ with $\sigma=2$, $\alpha=0.05$ two-sided, power $0.90$ ($\beta=0.10$, $z_\beta=1.282$).
  \item \textbf{Formula.} $n=\dfrac{(z_{\alpha/2}+z_\beta)^2\sigma^2}{\delta^2}$.
  \item \textbf{Compute.} $n=\dfrac{(1.96+1.282)^2(4)}{1}=42.0$.
  \item \textbf{Conclusion.} $n=42$.
\end{itemize}
\end{example}

\section{Chapter 10: Two Samples (book Chapter~4)}

\begin{formulabox}{Tests for a difference in means}
\textbf{$\sigma$'s known}
\[ z_0=\frac{\bar{x}_1-\bar{x}_2-\Delta_0}{\sqrt{\sigma_1^2/n_1+\sigma_2^2/n_2}} \]
\textbf{Pooled, equal variances ($n_1+n_2-2$ df)}
\[ \begin{aligned}t_0&=\frac{\bar{x}_1-\bar{x}_2-\Delta_0}{s_p\sqrt{1/n_1+1/n_2}}\\[4pt] s_p^2&=\frac{(n_1-1)s_1^2+(n_2-1)s_2^2}{n_1+n_2-2}\end{aligned} \]
\textbf{Welch, unequal variances}
\[ t_0^*=\frac{\bar{x}_1-\bar{x}_2-\Delta_0}{\sqrt{s_1^2/n_1+s_2^2/n_2}} \]
\textbf{Welch degrees of freedom (round down)}
\[ \nu=\frac{(a_1+a_2)^2}{\dfrac{a_1^2}{n_1-1}+\dfrac{a_2^2}{n_2-1}} \]
\[ a_i=\frac{s_i^2}{n_i} \]
\textbf{Paired ($n-1$ df)}
\[ t_0=\frac{\bar{d}-\Delta_0}{s_d/\sqrt n} \]
\end{formulabox}

\begin{formulabox}{Variances and proportions}
\textbf{Variance ratio}
\[ f_0=\frac{s_1^2}{s_2^2}\sim F_{n_1-1,\,n_2-1} \]
\[ f_{1-\alpha,u,v}=\frac1{f_{\alpha,v,u}} \]
\textbf{Two proportions (pooled $\hat{p}$ for the test)}
\[ \begin{aligned}z_0&=\frac{\hat{p}_1-\hat{p}_2}{\sqrt{\hat{p}(1-\hat{p})\big(\frac1{n_1}+\frac1{n_2}\big)}}\\[4pt] \hat{p}&=\frac{x_1+x_2}{n_1+n_2}\end{aligned} \]
Two-sided regions: $\abs{z_0}>z_{\alpha/2}$, $\abs{t_0}>t_{\alpha/2,\cdot}$, and $f_0>f_{\alpha/2}$ or $f_0<f_{1-\alpha/2}$.
\end{formulabox}

\begin{formulabox}{Two-sided intervals}
\begin{itemize}
  \item \textbf{$\sigma$'s known:} $(\bar{x}_1-\bar{x}_2)\pm z_{\alpha/2}\sqrt{\sigma_1^2/n_1+\sigma_2^2/n_2}$
  \item \textbf{Pooled:} $(\bar{x}_1-\bar{x}_2)\pm t_{\alpha/2,\,n_1+n_2-2}\,s_p\sqrt{1/n_1+1/n_2}$
  \item \textbf{Welch:} $(\bar{x}_1-\bar{x}_2)\pm t_{\alpha/2,\,\nu}\sqrt{s_1^2/n_1+s_2^2/n_2}$
  \item \textbf{Paired:} $\bar{d}\pm t_{\alpha/2,\,n-1}\,s_d/\sqrt n$
\end{itemize}
\textbf{Variance ratio}
\[ \begin{aligned}\frac{\sigma_1^2}{\sigma_2^2}&\ \ge\ \frac{s_1^2}{s_2^2}f_{1-\alpha/2,\,n_2-1,\,n_1-1}\\[4pt] \frac{\sigma_1^2}{\sigma_2^2}&\ \le\ \frac{s_1^2}{s_2^2}f_{\alpha/2,\,n_2-1,\,n_1-1}\end{aligned} \]
\textbf{Two proportions (unpooled)}
\[ \begin{aligned}&(\hat{p}_1-\hat{p}_2)\pm z_{\alpha/2}\,\widehat{\se}\\ &\widehat{\se}=\sqrt{\tfrac{\hat{p}_1\hat{q}_1}{n_1}+\tfrac{\hat{p}_2\hat{q}_2}{n_2}}\end{aligned} \]
\end{formulabox}

\begin{formulabox}{One-sided tests ($H_1$ uses $>$ or $<$): reject $H_0$ if}
The sign in $H_1$ points to the tail: $>$ means the statistic must be large and positive; $<$ means it must be large and negative.
\textbf{Difference of means ($z$ or $t$)}
\[ \begin{aligned} H_1:\mu_1-\mu_2>\Delta_0&:\ z_0>z_\alpha\ \ \text{or}\ \ t_0>t_{\alpha,\,\cdot}\\ H_1:\mu_1-\mu_2<\Delta_0&:\ z_0<-z_\alpha\ \ \text{or}\ \ t_0<-t_{\alpha,\,\cdot}\end{aligned} \]
\textbf{Variance ratio ($F$)}
\[ \begin{aligned} H_1:\sigma_1^2>\sigma_2^2&:\ f_0>f_{\alpha,\,n_1-1,\,n_2-1}\\ H_1:\sigma_1^2<\sigma_2^2&:\ f_0<f_{1-\alpha,\,n_1-1,\,n_2-1}\end{aligned} \]
\textbf{Two proportions ($z$)}
\[ \begin{aligned} H_1:p_1>p_2&:\ z_0>z_\alpha\\ H_1:p_1<p_2&:\ z_0<-z_\alpha\end{aligned} \]
\end{formulabox}

\begin{formulabox}{One-sided bounds: use $\alpha$, keep one limit}
\begin{itemize}
  \item \textbf{Means, upper bound (here $\sigma_D=\sqrt{\sigma_1^2/n_1+\sigma_2^2/n_2}$; with unknown $\sigma$ use $t_{\alpha,\cdot}$ and $s$):} $\mu_1-\mu_2\le(\bar{x}_1-\bar{x}_2)+z_\alpha\,\sigma_D$
  \item \textbf{Means, lower bound:} $\mu_1-\mu_2\ge(\bar{x}_1-\bar{x}_2)-z_\alpha\,\sigma_D$
  \item \textbf{Paired, upper bound:} $\mu_D\le\bar{d}+t_{\alpha,\,n-1}\tfrac{s_d}{\sqrt n}$
  \item \textbf{Paired, lower bound:} $\mu_D\ge\bar{d}-t_{\alpha,\,n-1}\tfrac{s_d}{\sqrt n}$
\end{itemize}
\textbf{Variance ratio, upper bound}
\[ \frac{\sigma_1^2}{\sigma_2^2}\le\frac{s_1^2}{s_2^2}f_{\alpha,\,n_2-1,\,n_1-1} \]
\textbf{Variance ratio, lower bound}
\[ \frac{\sigma_1^2}{\sigma_2^2}\ge\frac{s_1^2}{s_2^2}f_{1-\alpha,\,n_2-1,\,n_1-1} \]
\end{formulabox}

\begin{formulabox}{Sample size}
\textbf{Test, $\sigma$'s known}
\[ n\approx\frac{(z_{\alpha/2}+z_\beta)^2(\sigma_1^2+\sigma_2^2)}{\delta^2} \]
\textbf{Interval}
\[ n=\Big(\frac{z_{\alpha/2}}{E}\Big)^2(\sigma_1^2+\sigma_2^2) \]
\textbf{OC chart parameter, and equivalent $n$ for unequal sizes}
\[ d=\frac{\abs{\delta}}{\sqrt{\sigma_1^2+\sigma_2^2}} \]
\[ n=\frac{\sigma_1^2+\sigma_2^2}{\sigma_1^2/n_1+\sigma_2^2/n_2} \]
\end{formulabox}

\begin{notebox}
\begin{itemize}
  \item Pooled versus Welch: with $n_1=n_2$ they coincide. Pooled misleads when variances differ \emph{and} $n_1\ne n_2$. Decide from the normal plots, not a preliminary $F$-test; very different slopes mean use Welch.
  \item Paired: $\Var(\bar{D})=2\sigma^2(1-\rho)/n$. Pairing helps when units are heterogeneous and $\rho$ is large; it costs $n-1$ degrees of freedom. Never run an independent test on paired data.
  \item The $F$-test is not robust to non-normality; its two-sided $P$-value adds the tails of $f_0$ and $1/f_0$.
  \item For $n_1,n_2>40$ use $s_i$ in the $z$-test. Test the proportion difference with the \emph{pooled} $\hat{p}$ but build the interval with the \emph{unpooled} standard error.
  \item Cause and effect needs randomisation; observational differences can be confounded (Simpson's paradox).
\end{itemize}
\end{notebox}

\subsubsection*{Worked numericals}

\begin{example}[Pooled two-sample $t$-test]
\begin{itemize}
  \item \textbf{Given.} Method A: $n_1=10$, $\bar{x}_1=48.2$, $s_1=3.1$. Method B: $n_2=12$, $\bar{x}_2=45.1$, $s_2=3.7$. $H_0:\mu_1=\mu_2$ against $H_1:\mu_1\ne\mu_2$, $\alpha=0.05$, equal variances.
  \item \textbf{Pooled variance.} $s_p^2=\dfrac{9(9.61)+11(13.69)}{20}=11.854$, so $s_p=3.443$.
  \item \textbf{Standard error.} $3.443\sqrt{1/10+1/12}=1.474$.
  \item \textbf{Statistic.} $t_0=\dfrac{3.1}{1.474}=2.10$ with $20$ df.
  \item \textbf{Critical value.} $t_{0.025,20}=2.086$. Since $2.10>2.086$, reject.
  \item \textbf{Interval.} $3.1\pm2.086(1.474)$, i.e.\ $(0.03,\ 6.17)$. It excludes $0$.
  \item \textbf{Conclusion.} The means differ (barely). The interval agrees with the test.
\end{itemize}
\end{example}

\begin{example}[Welch test, same data]
\begin{itemize}
  \item \textbf{Given.} Same samples, but do not assume equal variances. $a_1=s_1^2/n_1=0.961$, $a_2=s_2^2/n_2=1.141$.
  \item \textbf{Statistic.} $t_0^*=\dfrac{3.1}{\sqrt{0.961+1.141}}=2.14$.
  \item \textbf{Degrees of freedom.} $\nu=\dfrac{(2.102)^2}{0.961^2/9+1.141^2/11}=19.99$, round down to $19$.
  \item \textbf{Critical value.} $t_{0.025,19}=2.093$. Since $2.14>2.093$, reject.
  \item \textbf{Conclusion.} Same conclusion as the pooled test.
\end{itemize}
\end{example}

\begin{example}[Paired $t$-test]
\begin{itemize}
  \item \textbf{Given.} Eight differences: $2.1,\ 1.4,\ -0.3,\ 3.0,\ 1.8,\ 0.9,\ 2.5,\ 1.2$. $H_0:\mu_D=0$.
  \item \textbf{Summary.} $\bar{d}=1.575$, $s_d=1.025$, $n=8$.
  \item \textbf{Statistic.} $t_0=\dfrac{1.575}{1.025/\sqrt8}=4.35$ with $7$ df.
  \item \textbf{Critical value.} $t_{0.005,7}=3.499$. Since $4.35>3.499$, reject even at $\alpha=0.01$.
  \item \textbf{Conclusion.} The mean difference is not zero.
\end{itemize}
\end{example}

\begin{example}[$F$-test for two variances]
\begin{itemize}
  \item \textbf{Given.} $s_1^2=9.61$ ($n_1=10$), $s_2^2=13.69$ ($n_2=12$). $H_0:\sigma_1^2=\sigma_2^2$, $\alpha=0.05$ two-sided.
  \item \textbf{Statistic.} $f_0=9.61/13.69=0.702$ with $(9,\,11)$ df.
  \item \textbf{Upper point.} $f_{0.025,9,11}=3.59$.
  \item \textbf{Lower point.} $f_{0.975,9,11}=1/f_{0.025,11,9}=1/3.91=0.256$.
  \item \textbf{Decision.} $0.256<0.702<3.59$, so do not reject.
  \item \textbf{Conclusion.} No evidence that the variances differ.
\end{itemize}
\end{example}

\begin{example}[Two proportions]
\begin{itemize}
  \item \textbf{Given.} Line 1: $15$ defective of $300$ ($\hat{p}_1=0.05$). Line 2: $20$ of $250$ ($\hat{p}_2=0.08$). $H_0:p_1=p_2$, two-sided, $\alpha=0.05$.
  \item \textbf{Pooled proportion.} $\hat{p}=35/550=0.0636$.
  \item \textbf{Statistic.} $z_0=\dfrac{-0.03}{\sqrt{0.0636(0.9364)(1/300+1/250)}}=\dfrac{-0.03}{0.0209}=-1.44$.
  \item \textbf{Decision.} $\abs{-1.44}<1.96$ ($P=0.15$), so do not reject.
  \item \textbf{Interval (unpooled).} $-0.03\pm1.96(0.0212)$, i.e.\ $(-0.072,\ 0.012)$.
  \item \textbf{Conclusion.} No evidence of a difference. The test uses the pooled $\hat{p}$; the interval uses the unpooled standard error.
\end{itemize}
\end{example}

\section{Chapter 11: Simple Linear Regression and Correlation (book Chapter~5)}

\begin{formulabox}{Least squares}
\textbf{Sums}
\[ \begin{aligned}S_{xx}&=\sum(x_i-\bar{x})^2\\ S_{xy}&=\sum(x_i-\bar{x})(y_i-\bar{y})\end{aligned} \]
\textbf{Estimates}
\[ \hat\beta_1=\frac{S_{xy}}{S_{xx}} \]
\begin{itemize}
  \item $\hat\beta_0=\bar{y}-\hat\beta_1\bar{x}$
  \item \textbf{Error variance (two parameters estimated):} $SS_E=S_{yy}-\hat\beta_1S_{xy}$
\end{itemize}
\[ \hat\sigma^2=\frac{SS_E}{n-2} \]
\textbf{Standard errors}
\[ \begin{aligned}\se(\hat\beta_1)&=\sqrt{\frac{\hat\sigma^2}{S_{xx}}}\\[4pt] \se(\hat\beta_0)&=\sqrt{\hat\sigma^2\Big(\frac1n+\frac{\bar{x}^2}{S_{xx}}\Big)}\end{aligned} \]
\end{formulabox}

\begin{formulabox}{Tests}
\textbf{Slope $t$-test ($n-2$ df)}
\[ t_0=\frac{\hat\beta_1-\beta_{1,0}}{\se(\hat\beta_1)} \]
\textbf{ANOVA form of the test for $\beta_1=0$}
\[ \begin{aligned}F_0&=\frac{MS_R}{MS_E}=t_0^2\\ &(1,\ n-2\text{ df})\\ SS_T&=SS_R+SS_E\\ SS_R&=\hat\beta_1S_{xy}\end{aligned} \]
\textbf{Coefficient of determination}
\[ R^2=\frac{SS_R}{SS_T} \]
\end{formulabox}

\begin{formulabox}{Intervals ($t=t_{\alpha/2,\,n-2}$)}
\begin{itemize}
  \item \textbf{Slope:} $\hat\beta_1\pm t\,\se(\hat\beta_1)$
\end{itemize}
\textbf{Mean response at $x_0$}
\[ \begin{aligned}&\hat{y}_0\pm t\,\se_{\mu}\\ &\se_{\mu}=\sqrt{\hat\sigma^2\Big[\frac1n+\frac{(x_0-\bar{x})^2}{S_{xx}}\Big]}\end{aligned} \]
\textbf{New observation at $x_0$}
\[ \begin{aligned}&\hat{y}_0\pm t\,\se_{\text{pred}}\\ &\se_{\text{pred}}=\sqrt{\hat\sigma^2\Big[1+\frac1n+\frac{(x_0-\bar{x})^2}{S_{xx}}\Big]}\end{aligned} \]
\end{formulabox}

\begin{formulabox}{Correlation}
\textbf{Sample correlation}
\[ r=\frac{S_{xy}}{\sqrt{S_{xx}S_{yy}}} \]
\[ \hat\beta_1=r\sqrt{\frac{S_{yy}}{S_{xx}}} \]
\begin{itemize}
  \item $r^2=R^2$
\end{itemize}
\textbf{Test of $\rho=0$ ($n-2$ df; same as the slope test)}
\[ t_0=\frac{r\sqrt{n-2}}{\sqrt{1-r^2}} \]
\textbf{Approximate interval for $\rho$ (Fisher)}
\[ \tanh\!\Big(\operatorname{arctanh}r\pm\frac{z_{\alpha/2}}{\sqrt{n-3}}\Big) \]
\end{formulabox}

\begin{notebox}
\begin{itemize}
  \item $\hat\sigma^2$ divides by $n-2$ (two parameters estimated). $\operatorname{cov}(\hat\beta_0,\hat\beta_1)=-\bar{x}\sigma^2/S_{xx}$. The $\hat\beta$'s are unbiased and BLUE.
  \item The $t$-test allows one-sided alternatives; the $F$-test does not. One-sided bounds use $t_{\alpha,n-2}$.
  \item Both intervals are narrowest at $x_0=\bar{x}$; the prediction interval is always wider. The confidence level applies to one $x_0$, not the whole line.
  \item $R^2$ does not measure the slope or whether the model is correct, and it always rises with more terms.
  \item Residuals: standardised $d_i=e_i/\hat\sigma$ (about $95\%$ within $\pm2$). A curve means the wrong form; a funnel means transform $y$ ($\sqrt{y}$, $\ln y$, $1/y$). Linearise $y=\beta_0e^{\beta_1x}$ by $\ln y$, $y=\beta_0x^{\beta_1}$ by logs of both.
  \item Correlation is not causation, and $\rho=0$ is not independence. Textbook values: oxygen purity $\hat\beta_1=14.947$, $\hat\beta_0=74.283$, $R^2=0.877$.
\end{itemize}
\end{notebox}

\subsubsection*{Worked numericals}

\begin{example}[Fit a regression line]
\begin{itemize}
  \item \textbf{Given.} $x=1,2,3,4,5$ and $y=2.1,\ 3.9,\ 6.2,\ 7.8,\ 10.1$. So $n=5$, $\bar{x}=3$, $\bar{y}=6.02$.
  \item \textbf{Sums.} $S_{xx}=10$, $S_{xy}=19.9$, $S_{yy}=39.708$.
  \item \textbf{Slope.} $\hat\beta_1=19.9/10=1.99$.
  \item \textbf{Intercept.} $\hat\beta_0=6.02-1.99(3)=0.05$.
  \item \textbf{Error variance.} $SS_E=39.708-1.99(19.9)=0.107$; $\hat\sigma^2=0.107/3=0.0357$ ($3$ df).
  \item \textbf{Conclusion.} Fitted line: $\hat{y}=0.05+1.99x$.
\end{itemize}
\end{example}

\begin{example}[Test the slope and report $R^2$]
\begin{itemize}
  \item \textbf{Given.} Continue from the fitted line. $H_0:\beta_1=0$, $\alpha=0.05$.
  \item \textbf{Standard error.} $\se(\hat\beta_1)=\sqrt{0.0357/10}=0.0597$.
  \item \textbf{Statistic.} $t_0=1.99/0.0597=33.3$.
  \item \textbf{Critical value.} $t_{0.025,3}=3.182$. Since $33.3>3.182$, reject.
  \item \textbf{$R^2$.} $1-0.107/39.708=0.997$.
  \item \textbf{Conclusion.} The slope is highly significant, and $99.7\%$ of the variability is explained.
\end{itemize}
\end{example}

\begin{example}[Intervals at $x_0=3.5$]
\begin{itemize}
  \item \textbf{Given.} Continue from the fitted line. $\hat{y}_0=0.05+1.99(3.5)=7.015$. Use $t_{0.025,3}=3.182$.
  \item \textbf{Mean response.} $\se=\sqrt{0.0357\big(\tfrac15+\tfrac{0.25}{10}\big)}=0.0896$, so $7.015\pm3.182(0.0896)=7.015\pm0.285$.
  \item \textbf{New observation.} $\se=\sqrt{0.0357\big(1+\tfrac15+\tfrac{0.25}{10}\big)}=0.209$, so $7.015\pm3.182(0.209)=7.015\pm0.665$.
  \item \textbf{Conclusion.} The prediction interval is wider because it includes the variation of a single new observation.
\end{itemize}
\end{example}

\section{Chapter 12: Multiple Linear Regression (book Chapter~6)}

\begin{formulabox}{Estimation}
\begin{itemize}
  \item \textbf{Model and estimator ($p=k+1$ parameters):} $\mathbf{y}=X\boldsymbol\beta+\boldsymbol\varepsilon$
  \item $\hat{\boldsymbol\beta}=(X\T X)^{-1}X\T\mathbf{y}$
  \item \textbf{Covariance and standard error:} $\Cov[\hat{\boldsymbol\beta}]=\sigma^2(X\T X)^{-1}$
  \item $\se(\hat\beta_j)=\sqrt{\hat\sigma^2C_{jj}}$
\end{itemize}
\textbf{Error variance}
\[ \hat\sigma^2=\frac{SS_E}{n-p} \]
\begin{itemize}
  \item \textbf{Hat matrix, fitted values, residuals:} $H=X(X\T X)^{-1}X\T$
  \item $\hat{\mathbf{y}}=H\mathbf{y}$
  \item $\mathbf{e}=(I-H)\mathbf{y}$
\end{itemize}
\end{formulabox}

\begin{formulabox}{Tests}
\textbf{Overall (upper tail)}
\[ \begin{aligned}F_0&=\frac{SS_R/k}{SS_E/(n-p)}\\ F_0&\sim F_{k,\,n-p}\end{aligned} \]
\textbf{One coefficient, given the others ($n-p$ df)}
\[ t_0=\frac{\hat\beta_j}{\se(\hat\beta_j)} \]
\textbf{A subset of $r$ coefficients (partial $F$; for $r=1$, $F_0=t_0^2$)}
\[ \begin{aligned}F_0&=\frac{SS_R(\boldsymbol\beta_2\mid\boldsymbol\beta_1)/r}{MS_E(\text{full})}\\ F_0&\sim F_{r,\,n-p}\end{aligned} \]
\textbf{Adjusted $R^2$}
\[ R^2_{\mathrm{adj}}=1-\frac{SS_E/(n-p)}{SS_T/(n-1)} \]
\end{formulabox}

\begin{formulabox}{Intervals ($t=t_{\alpha/2,\,n-p}$)}
\begin{itemize}
  \item \textbf{Coefficient:} $\hat\beta_j\pm t\sqrt{\hat\sigma^2C_{jj}}$
  \item \textbf{Mean response:} $\hat\mu_0\pm t\sqrt{\hat\sigma^2\,h_{00}}$
\end{itemize}
\textbf{New observation (adds 1 inside the root)}
\[ \begin{aligned}&\hat{y}_0\pm t\sqrt{\hat\sigma^2(1+h_{00})}\\ &h_{00}=\mathbf{x}_0\T(X\T X)^{-1}\mathbf{x}_0\end{aligned} \]
\end{formulabox}

\begin{formulabox}{Diagnostics}
\textbf{Studentised residual and Cook's distance}
\[ r_i=\frac{e_i}{\sqrt{\hat\sigma^2(1-h_{ii})}} \]
\[ D_i=\frac{r_i^2}{p}\,\frac{h_{ii}}{1-h_{ii}} \]
\textbf{PRESS and variance inflation factor}
\[ \begin{aligned}\mathrm{PRESS}&=\sum\Big(\frac{e_i}{1-h_{ii}}\Big)^2\\ \mathrm{VIF}_j&=\frac1{1-R_j^2}\end{aligned} \]
\end{formulabox}

\begin{notebox}
\begin{itemize}
  \item High leverage: $h_{ii}>2p/n$. Influential: $D_i>1$. Beware hidden extrapolation.
  \item Model building: all possible regressions (minimum $MS_E$, $C_p\approx p$, using the full-model $MS_E$); stepwise uses $f_{\mathrm{in}}\ge f_{\mathrm{out}}$ (default $\alpha=0.15$); forward only adds, backward only removes. Selection makes the reported $P$-values optimistic.
  \item Indicator variables: $m$ levels need $m-1$ indicators.
  \item Multicollinearity: $\mathrm{VIF}>10$ is a concern (some use $4$--$5$). Sign: significant overall $F$ but insignificant individual $t$'s. Fixes: more data, drop terms, ridge $(X\T X+\lambda I)^{-1}X\T\mathbf{y}$. Interpolation can stay accurate; extrapolation cannot.
  \item Textbook values (wire bond, $n=25$): $\hat\beta=(2.264,\,2.744,\,0.01253)$, $\hat\sigma^2=5.236$, $R^2=0.9811$, $R^2_{\mathrm{adj}}=0.979$.
\end{itemize}
\end{notebox}

\subsubsection*{Worked numericals}

\begin{example}[Overall $F$-test and adjusted $R^2$]
\begin{itemize}
  \item \textbf{Given.} $k=3$ regressors, $n=20$ (so $p=4$, error df $16$). $SS_T=500$, $SS_R=420$, $SS_E=80$.
  \item \textbf{Statistic.} $F_0=\dfrac{420/3}{80/16}=\dfrac{140}{5}=28$.
  \item \textbf{Critical value.} $f_{0.05,3,16}=3.24$. Since $28>3.24$, reject.
  \item \textbf{$R^2$.} $420/500=0.84$.
  \item \textbf{Adjusted $R^2$.} $1-\dfrac{80/16}{500/19}=1-\dfrac{5}{26.32}=0.810$.
  \item \textbf{Conclusion.} The regression is significant.
\end{itemize}
\end{example}

\begin{example}[Partial $F$ for one regressor]
\begin{itemize}
  \item \textbf{Given.} Remove one regressor. The reduced model has $SS_R=400$, so $SS_R(\beta_j\mid\text{others})=420-400=20$. Use $MS_E=5$ from the full model.
  \item \textbf{Statistic.} $F_0=\dfrac{20/1}{5}=4.0$, and $t_0=\sqrt{4.0}=2.0$.
  \item \textbf{Critical value.} $f_{0.05,1,16}=4.49$ (equivalently $t_{0.025,16}=2.120$). Since $4.0<4.49$, do not reject.
  \item \textbf{Conclusion.} This regressor is not significant. Dropping it is reasonable.
\end{itemize}
\end{example}

\begin{example}[Leverage check]
\begin{itemize}
  \item \textbf{Given.} $p=4$, $n=20$. One observation has $h_{ii}=0.55$.
  \item \textbf{Cutoff.} $2p/n=8/20=0.40$.
  \item \textbf{Compare.} $0.55>0.40$.
  \item \textbf{Conclusion.} This observation has high leverage. Check Cook's distance ($D_i>1$ means influential).
\end{itemize}
\end{example}

\section{Chapter 13: Single-Factor Experiments, ANOVA (book Chapter~7)}

\begin{formulabox}{One-way ANOVA}
\begin{itemize}
  \item \textbf{Model and degrees of freedom:} $Y_{ij}=\mu+\tau_i+\varepsilon_{ij}$
  \item $N-1=(a-1)+a(n-1)$
\end{itemize}
\textbf{Sums of squares}
\[ SS_T=\sum\sum y_{ij}^2-\frac{y_{\cdot\cdot}^2}{N} \]
\[ SS_{\text{Tr}}=\frac1n\sum y_{i\cdot}^2-\frac{y_{\cdot\cdot}^2}{N} \]
\begin{itemize}
  \item $SS_E=SS_T-SS_{\text{Tr}}$
\end{itemize}
\textbf{Test (upper tail only)}
\[ F_0=\frac{MS_{\text{Tr}}}{MS_E}\sim F_{a-1,\,N-a} \]
\end{formulabox}

\begin{formulabox}{Follow-up comparisons}
\begin{itemize}
  \item \textbf{Treatment mean:} $\bar{y}_{i\cdot}\pm t_{\alpha/2,\,N-a}\sqrt{MS_E/n}$
  \item \textbf{Difference of two means:} $\bar{y}_{i\cdot}-\bar{y}_{j\cdot}\pm t_{\alpha/2,\,N-a}\sqrt{2MS_E/n}$
\end{itemize}
\textbf{Fisher LSD: means differ if}
\[ \begin{aligned}\abs{\bar{y}_{i\cdot}-\bar{y}_{j\cdot}}&>\mathrm{LSD}\\ \mathrm{LSD}&=t_{\alpha/2,\,N-a}\sqrt{MS_E\,c_{ij}}\\ c_{ij}&=\frac1{n_i}+\frac1{n_j}\end{aligned} \]
\begin{itemize}
  \item \textbf{Tukey: means differ if:} $\abs{\bar{y}_{i\cdot}-\bar{y}_{j\cdot}}>q_\alpha(a,f)\sqrt{MS_E/n}$
\end{itemize}
\textbf{Contrast (1 df, $\sum c_i=0$)}
\[ SS_C=\frac{\big(\sum c_iy_{i\cdot}\big)^2}{n\sum c_i^2} \]
\end{formulabox}

\begin{formulabox}{Random effects and blocks}
\begin{itemize}
  \item \textbf{Variance components (a negative $\hat\sigma_\tau^2$ is treated as 0):} $\hat\sigma^2=MS_E$
\end{itemize}
\[ \hat\sigma_\tau^2=\frac{MS_{\text{Tr}}-MS_E}{n} \]
\begin{itemize}
  \item \textbf{Randomized complete block design:} $SS_T=SS_{\text{Tr}}+SS_{\text{Blocks}}+SS_E$
  \item $\text{error df}=(a-1)(b-1)$
\end{itemize}
\[ F_0=\frac{MS_{\text{Tr}}}{MS_E} \]
\end{formulabox}

\begin{notebox}
\begin{itemize}
  \item $F$ is upper-tail only because $\Exp[MS_{\text{Tr}}]=\sigma^2+n\sum\tau_i^2/(a-1)$. A balanced design is robust to unequal variances.
  \item LSD is liberal (experimentwise error exceeds $\alpha$); Tukey keeps the overall error at $\alpha$. Both are two-sided.
  \item Residual checks: normal plot, residuals versus fitted ($\ln y$ or $\sqrt{y}$ if the variance grows with the mean), and versus run order.
  \item In a random-effects model test $\sigma_\tau^2$, not individual $\tau_i$.
  \item RCBD: blocking costs $b-1$ degrees of freedom; when in doubt, block. The $F$ for blocks is not a reliable test.
  \item Sample size: $\Phi^2=n\sum\tau_i^2/(a\sigma^2)$ with OC curves; textbook $a=5$, $\sum\tau_i^2=5\sigma^2$, power $0.90$ gives $n=6$.
  \item Textbook values: hardwood $f_0=19.60$ ($f_{0.01,3,20}=4.94$), LSD $=3.07$; fabric RCBD $f_0=75.13$ ($f_{0.05,3,12}=5.95$), LSD $=0.39$.
\end{itemize}
\end{notebox}

\subsubsection*{Worked numericals}

\begin{example}[One-way ANOVA]
\begin{itemize}
  \item \textbf{Given.} Three machine settings, $4$ runs each. A: $12,14,13,15$ (mean $13.5$). B: $16,18,17,19$ (mean $17.5$). C: $13,15,14,16$ (mean $14.5$). Grand mean $15.17$, $N=12$, $a=3$.
  \item \textbf{Treatments.} $SS_{\text{Tr}}=4[(13.5-15.17)^2+(17.5-15.17)^2+(14.5-15.17)^2]=34.67$ ($2$ df).
  \item \textbf{Total.} $SS_T=49.67$ ($11$ df).
  \item \textbf{Error.} $SS_E=49.67-34.67=15.00$ ($9$ df).
  \item \textbf{Mean squares.} $MS_{\text{Tr}}=17.33$, $MS_E=1.667$.
  \item \textbf{Statistic.} $f_0=17.33/1.667=10.4$.
  \item \textbf{Critical value.} $f_{0.05,2,9}=4.26$. Since $10.4>4.26$, reject.
  \item \textbf{Conclusion.} The three settings do not all have the same mean.
\end{itemize}
\end{example}

\begin{example}[Which means differ? (Fisher LSD)]
\begin{itemize}
  \item \textbf{Given.} Continue from the ANOVA above: $MS_E=1.667$, $n=4$, $9$ error df.
  \item \textbf{LSD.} $t_{0.025,9}\sqrt{2MS_E/n}=2.262\sqrt{2(1.667)/4}=2.07$.
  \item \textbf{Compare.} $B-A=4.0>2.07$ and $B-C=3.0>2.07$ differ; $C-A=1.0<2.07$ does not.
  \item \textbf{Conclusion.} Setting B differs from both A and C; A and C cannot be separated.
\end{itemize}
\end{example}

\section{Chapter 14: Several Factors (book Chapter~8)}

\begin{formulabox}{Two-factor factorial}
\begin{itemize}
  \item \textbf{Model:} $Y_{ijk}=\mu+\tau_i+\beta_j+(\tau\beta)_{ij}+\varepsilon_{ijk}$
  \item \textbf{Degrees of freedom:} $A:\ a-1,\quad B:\ b-1,\quad AB:\ (a-1)(b-1),\quad E:\ ab(n-1),\quad T:\ abn-1$
\end{itemize}
\textbf{Sums of squares}
\[ \begin{aligned}SS_A&=\frac1{bn}\sum y_{i\cdot\cdot}^2-\frac{y_{\cdots}^2}{abn}\\[4pt] SS_{\text{cells}}&=\frac1n\sum\sum y_{ij\cdot}^2-\frac{y_{\cdots}^2}{abn}\\[4pt] SS_{AB}&=SS_{\text{cells}}-SS_A-SS_B\end{aligned} \]
Each effect is tested with $F=MS/MS_E$; test the interaction $AB$ first.
\end{formulabox}

\begin{formulabox}{$2^k$ designs}
\textbf{Effect, sum of squares, coefficient ($C$ = contrast from the sign table)}
\[ \text{effect}=\frac{C}{2^{k-1}n} \]
\[ SS=\frac{C^2}{2^kn} \]
\begin{itemize}
  \item $\text{coefficient}=\tfrac12\,\text{effect}$
  \item \textbf{$2^2$ effects:} $A=\tfrac1{2n}[ab+a-b-(1)],\quad B=\tfrac1{2n}[ab+b-a-(1)],\quad AB=\tfrac1{2n}[ab+(1)-a-b]$
\end{itemize}
\textbf{Curvature test with centre points ($1,\ n_C-1$ df)}
\[ \begin{aligned}SS_{\text{curv}}&=\frac{n_Fn_C(\bar{y}_F-\bar{y}_C)^2}{n_F+n_C}\\[4pt] F_0&=\frac{SS_{\text{curv}}}{MS_E}\end{aligned} \]
\end{formulabox}

\begin{notebox}
\begin{itemize}
  \item Main effect (two levels) is the average at high minus the average at low. The interaction $AB$ is the difference of the diagonal averages. A significant interaction means read the cell means, not the main effects.
  \item One observation per cell: assume no interaction and use $MS_{AB}$ as error. Three factors need $n\ge2$. With $A$ fixed and $B$ random, test $A$ against $MS_{AB}$ and $B$ against $MS_E$.
  \item Single replicate: normal plot of the effects, or pool high-order interactions as error. Dropping $h$ factors projects onto a $2^{k-h}$ with $n2^h$ replicates.
  \item Blocking: $L=\sum\alpha_ix_i\pmod2$; confound the highest-order interaction ($2^3$ with $ABC$: blocks $\{(1),ab,ac,bc\}$ and $\{a,b,c,abc\}$). $2^p$ blocks lose $2^p-1$ effects.
  \item Fractions $2^{k-p}$: aliases come from multiplying by the defining relation ($A^2=I$); each effect has $2^p-1$ aliases; resolution is the length of the shortest word. III: main effects aliased with two-factor interactions; IV: main effects clear, two-factor interactions aliased in groups; V: both clear.
  \item Examples: $2^{3-1}_{\mathrm{III}}$ ($I=ABC$), $2^{4-1}_{\mathrm{IV}}$ ($I=ABCD$), $2^{6-2}_{\mathrm{IV}}$ ($I=ABCE=BCDF=ADEF$), $2^{7-4}_{\mathrm{III}}$ saturated. The alternate fraction $I=-ABC$ de-aliases.
  \item Response surfaces: steepest ascent far from the optimum; central composite design and canonical analysis near it.
\end{itemize}
\end{notebox}

\subsubsection*{Worked numericals}

\begin{example}[Effects in a $2^2$ factorial]
\begin{itemize}
  \item \textbf{Given.} $n=3$ replicates. Treatment totals: $(1)=80$, $a=100$, $b=60$, $ab=90$.
  \item \textbf{Effect of $A$.} $\dfrac{1}{2n}[ab+a-b-(1)]=\dfrac{90+100-60-80}{6}=8.33$.
  \item \textbf{Effect of $B$.} $\dfrac{1}{6}[ab+b-a-(1)]=\dfrac{90+60-100-80}{6}=-5.00$.
  \item \textbf{Interaction $AB$.} $\dfrac{1}{6}[ab+(1)-a-b]=\dfrac{90+80-100-60}{6}=1.67$.
  \item \textbf{Sums of squares.} $SS=C^2/(4n)=C^2/12$: $SS_A=50^2/12=208.3$, $SS_B=30^2/12=75.0$, $SS_{AB}=10^2/12=8.33$.
  \item \textbf{Regression.} Coefficient $=$ effect$/2$: $\hat{y}=\bar{y}+4.17x_A-2.5x_B+0.83x_Ax_B$ (coded $\pm1$).
  \item \textbf{Conclusion.} The effect of $A$ is large and positive, $B$ is negative, and the interaction is small.
\end{itemize}
\end{example}

\begin{example}[Aliases in a fraction]
\begin{itemize}
  \item \textbf{Given.} $2^{3-1}$ design with defining relation $I=ABC$.
  \item \textbf{Multiply by $ABC$.} Use $A^2=B^2=C^2=I$: $A\cdot ABC=BC$, so $A=BC$; likewise $B=AC$ and $C=AB$.
  \item \textbf{Resolution.} The shortest word in $I=ABC$ has $3$ letters, so the design has resolution III.
  \item \textbf{Conclusion.} Each main effect is aliased with a two-factor interaction.
\end{itemize}
\end{example}

\section{Chapter 15: Nonparametric Statistics (book Chapter~9)}

\begin{formulabox}{Rank-based tests}
\textbf{Sign test: $R^+\sim\mathrm{Bin}(n,\tfrac12)$; reject if $\min(r^+,r^-)\le r^*$; for $n\ge10$}
\[ z_0=\frac{r^+-0.5n}{0.5\sqrt n} \]
\textbf{Signed-rank: $W^++W^-=\frac{n(n+1)}2$; reject if $\min(w^+,w^-)\le w^*$; for $n\ge20$}
\[ \mu_W=\frac{n(n+1)}{4} \]
\[ \sigma_W^2=\frac{n(n+1)(2n+1)}{24} \]
\textbf{Rank-sum: $W_1+W_2=\frac{(n_1+n_2)(n_1+n_2+1)}2$; for $n_1,n_2>8$}
\[ \begin{aligned}\mu_{W_1}&=\frac{n_1(n_1+n_2+1)}{2}\\[4pt] \sigma^2_{W_1}&=\frac{n_1n_2(n_1+n_2+1)}{12}\end{aligned} \]
\textbf{Kruskal--Wallis (reject for large $H$)}
\[ \begin{aligned}H&=\frac{12}{N(N+1)}\sum\frac{R_{i\cdot}^2}{n_i}\\ &\quad-3(N+1)\\ H&\approx\chi^2_{a-1}\end{aligned} \]
\end{formulabox}

\begin{notebox}
\begin{itemize}
  \item Sign test: drop ties. One-sided puts all of $\alpha$ in one tail, so use the table column for two-sided $2\alpha$; for $H_1:\tilde\mu>\tilde\mu_0$ reject when $r^-$ is small.
  \item Signed-rank (symmetric continuous population): rank $\abs{x_i-\mu_0}$ with average ranks for ties. One-sided: $H_1:\mu>\mu_0$ rejects if $w^-\le w^*$; $H_1:\mu<\mu_0$ rejects if $w^+\le w^*$. Normal approximation for $n\ge20$.
  \item Rank-sum (Mann--Whitney): $W_1$ is the rank sum of the smaller sample; reject if $w_1$ or $w_2\le w_\alpha$. Normal approximation for $n_1,n_2>8$.
  \item Efficiency versus the $t$-test: about $0.95$ for the Wilcoxon tests with normal data and never below $0.86$; the sign test is about $2/\pi\approx0.64$.
  \item Kruskal--Wallis: use the tie-corrected form with many ties. If ANOVA on ranks disagrees with ANOVA on the raw data, trust the ranks.
  \item Textbook values: sign test $n=20$, $r^+=14$, $P=0.1153$; signed-rank $w=60$ against $w^*=52$.
\end{itemize}
\end{notebox}

\subsubsection*{Worked numericals}

\begin{example}[Sign test]
\begin{itemize}
  \item \textbf{Given.} $H_0:\tilde\mu=\tilde\mu_0$ against $H_1:\tilde\mu\ne\tilde\mu_0$. Of $n=12$ nonzero differences, $9$ are positive. Under $H_0$, $R^+\sim\mathrm{Bin}(12,\tfrac12)$.
  \item \textbf{Tail probability.} $\Prob(R^+\ge9)=\dfrac{220+66+12+1}{4096}=0.0730$.
  \item \textbf{$P$-value.} $2(0.0730)=0.146$.
  \item \textbf{Decision.} $0.146>0.05$, so do not reject.
  \item \textbf{Table check.} $\min(9,3)=3$ and $r^*_{0.05}=2$ for $n=12$; $3>2$, so again do not reject.
  \item \textbf{Conclusion.} No evidence that the median differs from $\tilde\mu_0$.
\end{itemize}
\end{example}

\begin{example}[Wilcoxon signed-rank test]
\begin{itemize}
  \item \textbf{Given.} Differences: $1.2,\ -0.4,\ 2.0,\ 0.8,\ -0.6,\ 1.5,\ 0.9,\ 2.2$. $H_0:\mu_D=0$, two-sided, $\alpha=0.05$.
  \item \textbf{Ranks of $\abs{d}$.} $0.4\to1$, $0.6\to2$, $0.8\to3$, $0.9\to4$, $1.2\to5$, $1.5\to6$, $2.0\to7$, $2.2\to8$.
  \item \textbf{Rank sums.} $W^-=1+2=3$ and $W^+=36-3=33$, so $w=\min=3$.
  \item \textbf{Critical value.} $w^*_{0.05}=3$ for $n=8$. Since $3\le3$, reject.
  \item \textbf{Conclusion.} The mean difference is not zero.
\end{itemize}
\end{example}

\begin{example}[Large-sample rank-sum test]
\begin{itemize}
  \item \textbf{Given.} $n_1=n_2=10$, observed $w_1=85$. $H_0:\mu_1=\mu_2$, $\alpha=0.05$ two-sided.
  \item \textbf{Mean.} $\mu_{W_1}=\dfrac{10(21)}{2}=105$.
  \item \textbf{Standard deviation.} $\sigma_{W_1}=\sqrt{\dfrac{10\cdot10\cdot21}{12}}=13.23$.
  \item \textbf{Statistic.} $z_0=\dfrac{85-105}{13.23}=-1.51$.
  \item \textbf{Decision.} $\abs{-1.51}<1.96$, so do not reject.
  \item \textbf{Conclusion.} No evidence that the means differ.
\end{itemize}
\end{example}

\section{Chapter 16: Statistical Quality Control (book Chapter~10)}

\begin{formulabox}{Control limits}
\begin{itemize}
  \item \textbf{Shewhart chart ($\alpha=0.0027$ per point):} $\mu_W\pm3\sigma_W$
  \item \textbf{$\bar{X}$ and $R$ charts (or $\bar{\bar{x}}\pm A_3\bar{s}$ with $\hat\sigma=\bar{s}/c_4$):} $\bar{\bar{x}}\pm A_2\bar{r}$
\end{itemize}
\[ \hat\sigma=\frac{\bar{r}}{d_2} \]
\begin{itemize}
  \item $\mathrm{UCL}_R=D_4\bar{r},\quad \mathrm{LCL}_R=D_3\bar{r}$
\end{itemize}
\textbf{Individuals and moving range}
\[ \bar{x}\pm\frac{3\,\overline{mr}}{1.128} \]
\begin{itemize}
  \item $\mathrm{UCL}_{MR}=3.267\,\overline{mr}$
\end{itemize}
\textbf{$P$ chart (and $NP$)}
\[ \bar{p}\pm3\sqrt{\frac{\bar{p}(1-\bar{p})}{n}} \]
\begin{itemize}
  \item $n\bar{p}\pm3\sqrt{n\bar{p}(1-\bar{p})}$
\end{itemize}
\textbf{$U$ chart (and $C$)}
\[ \bar{u}\pm3\sqrt{\frac{\bar{u}}{n}} \]
\begin{itemize}
  \item $\bar{c}\pm3\sqrt{\bar{c}}$
\end{itemize}
\textbf{Constants}
\[ \begin{array}{c|cccc}n&2&3&4&5\\\hline A_2&1.880&1.023&0.729&0.577\\D_3&0&0&0&0\\D_4&3.267&2.575&2.282&2.114\\d_2&1.128&1.693&2.059&2.326\end{array} \]
\end{formulabox}

\begin{formulabox}{Capability, run length and CUSUM}
\textbf{Process capability}
\[ \begin{aligned}\mathrm{PCR}&=\frac{\mathrm{USL}-\mathrm{LSL}}{6\sigma}\\[4pt] \mathrm{PCR}_k&=\min(c_U,\ c_L)\\ c_U&=\frac{\mathrm{USL}-\mu}{3\sigma}\\ c_L&=\frac{\mu-\mathrm{LSL}}{3\sigma}\end{aligned} \]
\textbf{Average run length}
\[ \begin{aligned}\mathrm{ARL}&=\frac1p\\ \mathrm{ARL}_0&=\frac1{0.0027}=370\\ \mathrm{ARL}_{\text{out}}&=\frac1{1-\beta}\end{aligned} \]
\begin{itemize}
  \item \textbf{Tabular CUSUM (signal if either exceeds $H$):} $S_H(i)=\max[0,\,x_i-(\mu_0+K)+S_H(i-1)]$
  \item $S_L(i)=\max[0,\,(\mu_0-K)-x_i+S_L(i-1)]$
\end{itemize}
Design: $K=\tfrac12\sigma$ and $H=4\sigma$ or $5\sigma$; a $1\sigma$ shift is detected in $8.38$ or $10.4$ samples.
\end{formulabox}

\begin{notebox}
\begin{itemize}
  \item A control chart is a repeated hypothesis test. Never plot specification limits on it.
  \item Western Electric rules: (1) a point beyond $3\sigma$; (2) two of three beyond $2\sigma$ on one side; (3) four of five beyond $1\sigma$ on one side; (4) eight in a row on one side. Do not add them to a CUSUM.
  \item Build limits from $20$--$25$ preliminary samples of size $4$--$6$; remove samples with assignable causes and recompute; check the $R$ chart first.
  \item $\mathrm{PCR}=1$ means $0.27\%$ ($2700$ ppm) fallout; minimum $1.33$, critical $1.66$, six-sigma $2.0$ ($3.4$ ppm with a $1.5\sigma$ shift). $\mathrm{PCR}=\mathrm{PCR}_k$ iff centred. Capability is meaningful only for a stable process.
  \item A $1\sigma$ shift with $n=1$ takes $43.9$ points on average to detect. Narrower limits raise false alarms.
  \item Other tools: histogram, Pareto chart, fishbone diagram, defect-concentration diagram, scatter diagram, check sheet.
\end{itemize}
\end{notebox}

\subsubsection*{Worked numericals}

\begin{example}[$\bar{X}$ and $R$ chart limits]
\begin{itemize}
  \item \textbf{Given.} $20$ samples of size $n=5$: $\bar{\bar{x}}=33.3$, $\bar{r}=5.8$. Constants for $n=5$: $A_2=0.577$, $D_3=0$, $D_4=2.114$, $d_2=2.326$.
  \item \textbf{$\bar{X}$ chart.} $33.3\pm0.577(5.8)=33.3\pm3.35$, i.e.\ $(29.95,\ 36.65)$.
  \item \textbf{$R$ chart.} $\mathrm{UCL}=2.114(5.8)=12.26$, $\mathrm{LCL}=0$.
  \item \textbf{Process sigma.} $\hat\sigma=\bar{r}/d_2=5.8/2.326=2.49$.
  \item \textbf{Conclusion.} Plot each sample mean and range against these limits.
\end{itemize}
\end{example}

\begin{example}[$P$ chart limits]
\begin{itemize}
  \item \textbf{Given.} Average fraction defective $\bar{p}=0.40$ ($800$ defectives in $2000$ units); subgroup size $n=100$.
  \item \textbf{Standard error.} $\sqrt{\dfrac{0.40(0.60)}{100}}=0.049$.
  \item \textbf{Limits.} $0.40\pm3(0.049)=0.40\pm0.147$.
  \item \textbf{Conclusion.} $\mathrm{LCL}=0.253$, $\mathrm{UCL}=0.547$.
\end{itemize}
\end{example}

\begin{example}[Process capability]
\begin{itemize}
  \item \textbf{Given.} $\mathrm{LSL}=90$, $\mathrm{USL}=110$, process mean $\mu=102$, $\sigma=2$.
  \item \textbf{$\mathrm{PCR}$.} $\dfrac{110-90}{6(2)}=1.67$.
  \item \textbf{$\mathrm{PCR}_k$.} $\min\Big(\dfrac{110-102}{6},\ \dfrac{102-90}{6}\Big)=\min(1.33,\ 2.00)=1.33$.
  \item \textbf{Fallout.} Mainly above the USL: $\Prob(X>110)=\Prob(Z>4)\approx3\times10^{-5}$.
  \item \textbf{Conclusion.} The process is capable but off-centre, because $\mathrm{PCR}_k<\mathrm{PCR}$.
\end{itemize}
\end{example}

\begin{example}[Average run length]
\begin{itemize}
  \item \textbf{Given.} A $3\sigma$ $\bar{X}$ chart.
  \item \textbf{In control.} $\mathrm{ARL}_0=1/0.0027=370$ samples between false alarms.
  \item \textbf{After a shift.} If each sample signals with probability $p=0.50$, then $\mathrm{ARL}=1/0.50=2$ samples.
  \item \textbf{Conclusion.} A larger sample size or more frequent sampling shortens the detection time.
\end{itemize}
\end{example}

\begin{example}[Tabular CUSUM]
\begin{itemize}
  \item \textbf{Given.} Target $\mu_0=99$, $K=1$, $H=10$. Observations $x=100,\ 102,\ 101,\ 103$. Upper sum: $S_H(i)=\max[0,\ x_i-(\mu_0+K)+S_H(i-1)]$, start at $0$.
  \item \textbf{$i=1$.} $S_H=\max(0,\ 100-100+0)=0$.
  \item \textbf{$i=2$.} $S_H=\max(0,\ 102-100+0)=2$.
  \item \textbf{$i=3$.} $S_H=\max(0,\ 101-100+2)=3$.
  \item \textbf{$i=4$.} $S_H=\max(0,\ 103-100+3)=6$.
  \item \textbf{Conclusion.} All values are below $H=10$, so there is no signal yet, but the sum keeps rising, which suggests an upward drift.
\end{itemize}
\end{example}

\section{Which procedure? A quick decision guide}

\begin{enumerate}
  \item \emph{One mean?} $\sigma$ known (or $n\ge40$): $z$. $\sigma$ unknown, normal: $t$. Not normal and small: sign or signed-rank.
  \item \emph{One variance?} $\chi^2$. \emph{One proportion?} $z$ ($np,nq\ge5$) or exact binomial.
  \item \emph{Two means, independent?} Known $\sigma$: $z$. Unknown equal: pooled $t$. Unknown unequal (or in doubt): Welch. Not normal: rank-sum.
  \item \emph{Two means, matched pairs?} Paired $t$ on the differences (or sign / signed-rank).
  \item \emph{Two variances?} $F$ (not robust). \emph{Two proportions?} Pooled $z$ for the test, unpooled for the interval.
  \item \emph{More than two means?} One factor: ANOVA $F$, then LSD or Tukey; non-normal: Kruskal--Wallis. With a nuisance factor: RCBD. Several factors: factorial ANOVA, $2^k$, fractions.
  \item \emph{Relationship?} One regressor: simple regression, $t$ or $F$ for $\beta_1$; several: multiple regression with partial $F$; both variables random: correlation.
  \item \emph{Distribution fit or independence of categories?} $\chi^2$ goodness of fit or contingency table.
  \item \emph{Monitoring a process?} Mean and spread: $\bar{X}$--$R$ or $\bar{X}$--$S$; individuals; attributes: $P$, $U$; small shifts: CUSUM.
\end{enumerate}

\begin{summarybox}
\begin{itemize}
  \item Count tails first: $\alpha/2$ for two-sided, $\alpha$ for one-sided, and one limit only for a bound.
  \item State the model assumptions and check them (normal plot, residuals, equal variances).
  \item Degrees of freedom: $n-1$ ($t$, $\chi^2$, paired), $n_1+n_2-2$ (pooled), $n-2$ (simple regression), $n-p$ (multiple), $N-a$ (one-way), $ab(n-1)$ (two-way), $(a-1)(b-1)$ (RCBD error), $k-p-1$ (goodness of fit).
  \item Round sample sizes \emph{up}; keep the strong claim in $H_1$; report an interval with every test.
\end{itemize}
\end{summarybox}
